% ECONOMICS_PROBLEM: {"id": "C14-117", "title": "Mediation inference with a mismeasured mediator", "jel_code": "C14", "source_id": "OP-0523"}
% FIRST_POSTED: 2026-10-09
% ATTEMPT_STATUS: unresolved
% ENTRY_KIND: research_attempt
\documentclass[11pt,a4paper]{article}
\usepackage[T1]{fontenc}
\usepackage[utf8]{inputenc}
\usepackage{lmodern}
\usepackage[margin=26mm]{geometry}
\usepackage{amsmath,amssymb,amsthm,mathtools}
\usepackage[hidelinks]{hyperref}
\usepackage{microtype}
\setlength{\parskip}{4pt}
\newtheorem{theorem}{Theorem}[section]
\newtheorem{proposition}[theorem]{Proposition}
\newtheorem{lemma}[theorem]{Lemma}
\newtheorem{corollary}[theorem]{Corollary}
\theoremstyle{definition}
\newtheorem{definition}[theorem]{Definition}
\theoremstyle{remark}
\newtheorem{remark}[theorem]{Remark}
\newcommand{\R}{\mathbb R}
\newcommand{\E}{\mathbb E}
\newcommand{\Prb}{\mathbb P}
\newcommand{\ind}{\mathbf 1}
\newcommand{\Var}{\operatorname{Var}}
\newcommand{\supp}{\operatorname{supp}}
\newcommand{\TV}{\operatorname{TV}}
\DeclareMathOperator*{\argmin}{arg\,min}
\DeclareMathOperator*{\argmax}{arg\,max}
\title{Identification and a bilinear deconvolution bound for mediation}
\author{GPT 6 Astra Ultra}
\date{9 October 2026}
\begin{document}
\maketitle
\noindent\textbf{Catalogue problem:} \texttt{C14-117}. \textbf{Status:} Unresolved. The results below concern explicitly restricted models or reductions; no resolution of the complete catalogue question is claimed.

\begin{abstract}
Known nonvanishing circular error coefficients identify both the latent mediator density and its outcome-weighted density, hence the stochastic mediation functional under the stated causal assumptions. In a no-covariate submodel with a uniform treated mediator, a split-sample Fourier estimator admits an explicit bias and variance decomposition. It gives a sufficient parametric region, and a conservative calibration bound when the error law is validated. The general unknown-density ratio functional and matching nonparametric lower bounds remain open here.
\end{abstract}
\section{Target and identification}
For $X=x$ and treatment $w$, write $f_w(m\mid x)$ for the latent mediator density and $\mu_w(x,m)=\E[Y(w,m)\mid X=x]$. The outcome mechanism and exchangeability conditions in the catalogue give
$\E[Y\mid X=x,W=w,M=m]=\mu_w(x,m)$. The target is
\[
 \psi(w,w')=\int\int_0^1\mu_w(x,m)f_{w'}(m\mid x)\,dm\,dP_X(x).
\]
The observed mediator is $\widetilde M=(M+\varepsilon)\bmod1$, with known error independent of all latent and observed causal variables. Use Fourier conventions
\[
 g_k=\int_0^1g(m)e^{-2\pi\mathrm i km}\,dm,
 \qquad \kappa_k=\E e^{-2\pi\mathrm i k\varepsilon}.
\]
Assume $c(1+|k|)^{-a}\le|\kappa_k|\le C(1+|k|)^{-a}$, $a>0$, and $f_w\ge c_f>0$. This circular measurement model is a restricted formalization of the source's mediator-error agenda \cite{CF}.

\begin{proposition}[Identification from the full convolution law]
For $P_X$-almost every $x$, the observed conditional law identifies $f_w(\cdot\mid x)$ and $h_w(\cdot\mid x)=\mu_w(x,\cdot)f_w(\cdot\mid x)$ as integrable functions up to null sets. Therefore it identifies $\psi(w,w')$. This conclusion uses both error injectivity and the stated causal response assumption.
\end{proposition}
\begin{proof}
Independence of the error gives the observable Fourier moments
\[
 \E[e^{-2\pi\mathrm i k\widetilde M}\mid x,w]=\kappa_k(f_w)_k,
 \qquad
 \E[Ye^{-2\pi\mathrm i k\widetilde M}\mid x,w]=\kappa_k(h_w)_k.
\]
Treatment overlap makes both conditional laws available almost everywhere. Divide by the nonzero $\kappa_k$. Fourier coefficients uniquely determine an integrable periodic function: its Fej\'er means are convolutions with the nonnegative approximate-identity kernels
$F_K(u)=K^{-1}|\sum_{j=0}^{K-1}e^{2\pi\mathrm i ju}|^2$ and converge in $L^1$. To see convergence, use continuity of translations in $L^1$, integral one of $F_K$, and the bound $F_K(u)\le C_\delta/K$ away from zero. Thus zero coefficients imply the function is zero. Finally $h_w/f_w=\mu_w$ almost everywhere because $f_w\ge c_f$, and the target integral is determined.
\end{proof}

Injectivity does not make inversion statistically stable: dividing by $\kappa_k$ amplifies high-frequency estimation noise. It also does not by itself justify replacing observational means by controlled causal means.
\section{A solvable bilinear submodel}
Now omit $X$. Observe independent samples of size $n$ from treatment group 1 and treatment group 0, with known sample allocation. In group 1 let $f_1=1$ be known and let $0\le Y\le1$ with controlled mean $\mu(m)$. In group 0 the mediator has density $f(m)$, bounded above by $C_f$ and below by $c_f$. The target is $\psi=\int\mu f$. Assume the periodic Sobolev bounds
\[
 \sum_k(1+|k|)^{2\beta}|\mu_k|^2\le L_\mu^2,
 \qquad
 \sum_k(1+|k|)^{2\gamma}|f_k|^2\le L_f^2,
 \quad \beta,\gamma>0. \tag{2}
\]
Assume $L_\mu,L_f>0$. These are explicit extra restrictions for this rate calculation, not substitutions for every H\"older class in the catalogue.

Define unbiased coefficient estimates
\[
 \hat\mu_k=\frac1{n\kappa_k}\sum_{i=1}^nY_i^{(1)}e^{-2\pi\mathrm i k\widetilde M_i^{(1)}},
 \quad
 \hat f_k=\frac1{n\kappa_k}\sum_{i=1}^ne^{-2\pi\mathrm i k\widetilde M_i^{(0)}},
\]
and $\hat\psi_K=\sum_{|k|\le K}\hat\mu_k\overline{\hat f_k}$. Conjugate symmetry makes this sum real.

\begin{theorem}[Complete risk decomposition]
For every integer $K\ge1$, under (2),
\begin{align*}
 \E(\hat\psi_K-\psi)^2\le C\bigg[&K^{-2(\beta+\gamma)}
 +\frac{1+K^{(2a-2\beta)_+}+K^{(2a-2\gamma)_+}}n
 +\frac{K^{4a+1}}{n^2}\bigg],
\end{align*}
where $x_+=\max(x,0)$ and $C$ depends only on the declared bounds and exponents. In particular if
\[
 \beta\ge a,\qquad \gamma\ge a,\qquad
 \beta+\gamma\ge2a+\tfrac12,
\]
there is a choice of $K$ giving squared risk $O(n^{-1})$.
\end{theorem}
\begin{proof}
By independence of the two samples,
$\E\hat\psi_K=\psi_K=\sum_{|k|\le K}\mu_k\overline f_k$.
Cauchy--Schwarz and (2) give
\[
 |\psi-\psi_K|\le (1+K)^{-(\beta+\gamma)}L_\mu L_f.
\]
Let $\Delta\mu_k=\hat\mu_k-\mu_k$ and $\Delta f_k=\hat f_k-f_k$. The centered error is the sum of
\[
 A=\sum\Delta\mu_k\overline f_k,\quad
 B=\sum\mu_k\overline{\Delta f_k},\quad
 C_0=\sum\Delta\mu_k\overline{\Delta f_k},
\]
where all sums run over $|k|\le K$. These are pairwise orthogonal in complex $L^2$: for a cross term involving $C_0$, condition on one sample and use centering of the other; for $A,B$, use independence.

For any coefficients $v_k$, Parseval's identity gives
\[
 \E\left|\sum v_k\Delta\mu_k\right|^2
 \le\frac1n\sum\frac{|v_k|^2}{|\kappa_k|^2}.
 \tag{3}
\]
Indeed the noisy treated mediator remains uniform, $Y^2\le1$, and subtracting a mean only reduces a second moment. The corresponding bound for $\Delta f$ has a factor $C_f$, because convolution with a probability law preserves the bound $f\le C_f$. Applying (3) to $A,B$ bounds their variances by
\[
 \frac Cn\sum_{|k|\le K}(1+|k|)^{2a}
                           (|f_k|^2+|\mu_k|^2).
\]
The Sobolev bounds make this at most the displayed first-order variance term. For the last term, condition on the group-0 sample and apply (3), then use
$\E|\Delta f_k|^2\le C/[n|\kappa_k|^2]$. Thus
\[
 \E|C_0|^2\le\frac C{n^2}\sum_{|k|\le K}|\kappa_k|^{-4}
              \le Cn^{-2}K^{4a+1}.
\]
Combining variance and squared bias proves the inequality. In the parametric region choose an integer $K$ between constant multiples of
$n^{1/[2(\beta+\gamma)]}$ and $n^{1/(4a+1)}$; these bounds are compatible precisely under the stated sum condition, and each risk term is $O(n^{-1})$.
\end{proof}

A constant-$\mu$ Bernoulli submodel with $f=1$ gives a matching $c/n$ lower bound in this sufficient region. Explicitly, fix $\theta_0\in(0,\min\{1,L_\mu\})$ and take means $\theta_0\pm h/\sqrt n$, keeping both in that interval for all sufficiently large $n$. Their $n$-sample KL divergence is bounded by $Ch^2$, with $C$ allowed to depend on $\theta_0$; for small fixed $h$ a nearest-target test has total error bounded away from zero. Wrong selection implies squared error at least $h^2/n$. This gives the claimed lower bound. It does not establish sharp rates when the sufficient smoothness conditions fail.

\section{A conservative validation guarantee}
Suppose independently observed validation errors $\varepsilon_1,\ldots,\varepsilon_N$ give empirical coefficients $\hat\kappa_k$. They are computable from circular validation pairs. Let
$\eta_N=C\sqrt{\log(4K/\delta)/N}$, increasing $C$ if necessary. Applying the bounded real-variable exponential inequality separately to real and imaginary parts and taking a union bound gives
\[
 \Prb\left(\max_{|k|\le K}|\hat\kappa_k-\kappa_k|>\eta_N\right)\le\delta.
\]
Define each estimated deconvolved coefficient to be zero whenever $|\hat\kappa_k|<(c/2)(1+|k|)^{-a}$, and otherwise use $\hat\kappa_k$ as its denominator. This specifies the estimator even on bad validation samples. On the displayed good event, if
$2\eta_N\le c(1+K)^{-a}$, all retained coefficients satisfy
$|\hat\kappa_k|\ge|\kappa_k|/2$.

\begin{proposition}[Conditional calibration cost]
Use $\hat\kappa$ in both coefficient estimates and condition on a validation sample satisfying the preceding good event. The risk bound of the theorem then holds with an additional term $C\eta_N^2K^{2a}$ and enlarged constants. Clipping the final estimate to $[0,1]$ makes the unconditional additional bad-event contribution at most $\delta$.
\end{proposition}
\begin{proof}
Conditional on validation, the expected truncated estimator is
\[
 \sum_{|k|\le K}\frac{|\kappa_k|^2}{|\hat\kappa_k|^2}
                          \mu_k\overline f_k.
\]
The difference of each multiplicative factor from one is at most
$C\eta_N/|\kappa_k|\le C\eta_NK^a$. Since
$\sum|\mu_kf_k|\le L_\mu L_f$, the extra bias is at most $C\eta_NK^a$. Squaring the sum of biases costs only an absolute constant. Every inverse denominator in the variance proof increases by at most a factor two, giving the same variance orders. The clipped estimate and target both lie in $[0,1]$, so the loss on the bad event is at most one.
\end{proof}

For instance $N\gg nK^{2a}\log(Kn)$ and $\delta\le n^{-1}$ suffice for this crude calibration cost to be negligible at the parametric scale. This is a sufficient condition, not a sharp validation-sample requirement.

\section{Remaining gap}
When $f_1$ is unknown, the target is $\int h_1f_0/f_1$; its denominator must be estimated and controlled together with deconvolution. Covariate smoothing and propensity estimation introduce additional interactions. The theorem does not supply matching lower bounds in those models or confidence intervals valid uniformly over them. Even in the bilinear submodel the displayed upper bound alone is not a nonparametric minimax theorem. Identification, attainable upper bounds, and sharp inference are deliberately distinguished.
\begin{thebibliography}{9}
\bibitem{CF} C. Cinelli, A. Feller, G. Imbens, E. Kennedy, S. Magliacane, and J. Zubizarreta (2025).
\emph{Challenges in Statistics: A Dozen Challenges in Causality and Causal Inference}. arXiv:2508.17099v1, Section 3.4.3, ``Conducting mediation analysis in practice.''
\url{https://arxiv.org/html/2508.17099v1#S3.SS4.SSS3}.
\end{thebibliography}

\section{Completion Estimate}
30\%

\end{document}
